\chapter{Contexte et Problématique}

\section{Introduction au Domaine}
L'éducation numérique a franchi un cap décisif au cours de la dernière décennie. Les institutions académiques s'appuient massivement sur les plateformes e-learning pour diffuser le savoir à une audience globale. Cependant, cette numérisation a révélé des vulnérabilités critiques, notamment en matière de suivi personnalisé de l'étudiant et de garantie de l'intégrité des examens passés à distance.

\section{La Problématique de l'Intégrité Académique}
Dans les environnements d'apprentissage virtuels, l'authentification continue de l'apprenant et la vérification de son honnêteté intellectuelle lors des épreuves d'évaluation posent un défi majeur. Les méthodes traditionnelles de surveillance humaine via webcam s'avèrent non seulement coûteuses en ressources, mais également invasives pour la vie privée et sujettes à l'erreur humaine liée à la fatigue visuelle des surveillants. 

\begin{infobox}[Le Défi de la Télésurveillance (Proctoring)]
Les systèmes actuels manquent souvent de finesse algorithmique. Ils génèrent un nombre disproportionné de fausses alertes (faux positifs), créant un environnement anxiogène pour l'étudiant, ou peinent à corréler des événements subtils (changement d'onglet combiné à un regard fuyant) pour détecter une véritable fraude.
\end{infobox}

\section{L'Émergence de l'Intelligence Artificielle Agentique}
Face à ces limites, l'intelligence artificielle agentique offre un paradigme nouveau. Un agent IA ne se limite pas à répondre à des requêtes ; il observe son environnement, analyse un contexte complexe et prend des décisions de manière autonome. L'application de ce paradigme au domaine de l'éducation (Agentic LMS) permet d'envisager des tuteurs virtuels infatigables, des générateurs d'évaluations adaptatives et des surveillants algorithmiques impartiaux.

\section{Objectifs du Projet Fraudly}
Le projet Fraudly vise à développer une plateforme éducative exhaustive intégrant organiquement cette intelligence agentique. L'objectif est double : d'une part, sublimer l'expérience d'apprentissage grâce à une personnalisation poussée par l'IA ; d'autre part, restaurer la confiance dans l'évaluation à distance grâce à un système de proctoring multimodal combinant vision par ordinateur et analyse comportementale.
